\section{总结与延伸}

这堂课的主线不是某个单一模型赢了，而是 native multimodal design 从“统一一切”逐步走向“统一接口、专用机制”。Chameleon 用离散 token 证明 early fusion 可行；Transfusion 保留共享 Transformer，却让 text 与 image 使用不同 objective；MoT 再按 modality 稀疏化参数；BAGEL 与 $\pi_{0.7}$ 把 specialization 推向统一内容模型和机器人控制。后半段则提醒：architecture 共存不等于 capability 正迁移，更不等于 physical intelligence。

\subsection{讲座结论与系统约束}

前面已经分别讨论 representation、objective、routing 与 transfer；本节回到 slide 056，把这些局部选择压缩成课程级系统结论。读图时应把 compute pressure 与 specialization 放在同一条因果链中。

\lecturefigure{slide-056.jpg}{课堂结论：open problems、physical-world gap、compute 与 specialization}{官方 deck 第 56 页；课堂 00:40:06--00:41:36}

\begin{knowledgebox}{读图：四点结论形成一条工程链}
第一，多模态仍是活跃且未定型的研究领域；第二，数字信息处理强，但物理世界能力远未完成；第三，image/video/audio 增加训练与 inference 基础设施压力；第四，短期内可能出现更多 capability-specific models。最后的问题不是“是否统一”，而是怎样在可负担计算与可靠交互下把专用模块组织成 coherent system。
\end{knowledgebox}

\teachervoice{00:40:41--00:41:32，老师把 computationally heavy 与 capability specialization 放在同一结论中：由于问题空间复杂、成本高，短期内专用模型增多是合理路径；如何再统一成系统将成为后续研究。}

\begin{importantbox}{本讲十条核心结论}
\begin{enumerate}
  \item Multimodal input、multimodal output 与 physical intelligence 是不同能力层。
  \item Token 是序列单元，不必是离散 vocabulary ID。
  \item Global autoregression 可以与 block 内 diffusion 共存。
  \item Chameleon 证明 discrete early fusion 可行，也暴露 quantization bottleneck。
  \item Transfusion 说明共享 backbone 不要求共享 loss。
  \item Modality gap 是 specialization 动机，不是因果证明。
  \item MoT 用已知 modality deterministic routing，MoE 用 learned gating。
  \item Sparse model 必须同时报告 total 与 active parameters 及 matched compute。
  \item Understanding 与 generation 的 transfer 必须按方向和 baseline 实验验证。
  \item 当前 language 是高效 reasoning scaffold，但 full physical-world intelligence 仍需闭环状态、控制与安全。
\end{enumerate}
\end{importantbox}

\subsection{一张可执行的 architecture review checklist}

\begin{knowledgebox}{从概念演示到可部署系统}
\begin{enumerate}
  \item \textbf{Representation}：每个 modality 是离散 code、continuous semantic embedding 还是 generative latent？
  \item \textbf{Ordering}：跨 modality 的 causal order 与 boundary 怎样定义？
  \item \textbf{Objective}：每种输出用什么 loss，怎样 normalization 与 weighting？
  \item \textbf{Routing}：哪些层共享，哪些按 modality 或 learned expert 稀疏化？
  \item \textbf{Accounting}：total/active parameters、tokens、FLOPs、sampling steps 与 wall-clock 是否同时报告？
  \item \textbf{Transfer}：联合训练相对 matched single-task baseline 是否真正改善？
  \item \textbf{Serving}：streaming、latency、KV/cache、non-text decoder 与并发怎样预算？
  \item \textbf{Embodiment}：若进入物理世界，是否有 state estimation、recovery、safety 与闭环 success metric？
\end{enumerate}
\end{knowledgebox}

\subsection{Deck appendix 的处理边界}

官方 PDF 还有三页 Interaction Models（pages 53--55），但完整录像从 broader multimodal intelligence 的 page 52 直接跳到 conclusion page 56，主讲与 Q\&A 均未展示或讲解这三页。按 source-first 原则，它们保留在 \texttt{slides-images/} 和 selection table 中，但标记 optional，不作为课堂内容重建。这样既保留官方 deck 完整性，也不把未讲内容误写成 teacher voice。

\begin{warningbox}{不能从本讲推出的结论}
不能推出所有 modality 最终都会共享一个 representation；不能推出 MoT 对 image understanding 同样有效；不能仅凭更低 diffusion loss 宣称生成更安全或更真实；不能把 BAGEL/$\pi_{0.7}$ 当作 full physical intelligence；也不能把 language 当前的 reasoning 优势解释为任何智能都必须先翻译成文本。
\end{warningbox}

\subsection{拓展阅读与 lecture snapshot}

建议按设计变量阅读一手论文：先用 Chameleon v2 理解 discrete early fusion，再用 Transfusion v1 对照 continuous diffusion objective；随后读 Mind the Gap v2 与 MoT v2，区分 representation geometry 与 architecture evidence；再读 LMFusion v4 看冻结 language path 的扩展方式，最后用 BAGEL v3 与 $\pi_{0.7}$ v2 比较 digital unified generation 与 embodied steering。

\begin{knowledgebox}{Lecture snapshot：2026-05-21}
本讲义冻结到课堂日前公开版本：
\begin{itemize}
  \item Chameleon \texttt{2405.09818v2}、VQ-VAE-2 \texttt{1906.00446v1}；
  \item Transfusion \texttt{2408.11039v1}、Mind the Gap \texttt{2203.02053v2}；
  \item MoT \texttt{2411.04996v2}、LMFusion \texttt{2412.15188v4}；
  \item BAGEL \texttt{2505.14683v3}、$\pi_{0.7}$ \texttt{2604.15483v2}。
\end{itemize}
所有版本均早于 2026-05-21；课后产品变化不倒灌为课堂事实。
\end{knowledgebox}

最终，native multimodal intelligence 不是把更多输入接到 LLM，也不是追求一张“统一架构图”。它是一组可审计的取舍：抽象与细节、共享与专用、理解与生成、数字内容与物理行动、模型质量与系统成本。真正成熟的系统会保留统一控制面，同时承认不同 modality 需要不同表示、目标、容量和验证方式。

\end{document}
